\section{总结与延伸}

\subsection{本章小结}

\begin{figure}[H]
\centering
\includegraphics[width=0.95\textwidth]{slides-images/slide-107.jpg}
\caption{对比学习部分的课程总结\protect\footnotemark}
\end{figure}
\footnotetext{Slides 第107页。}

\begin{importantbox}{本讲核心结论}
\begin{enumerate}
    \item 预文本任务利用数据本身自动生成监督信号，能在没有人工标签时学到有用表示。
    \item 图像变换类任务关注视觉常识，包括旋转、拼图、修补、颜色化和 MAE。
    \item 对比学习把问题抽象成“拉近正样本、推远负样本”，InfoNCE 是核心目标。
    \item SimCLR、MoCo、CPC、DINO 分别代表了不同的实现路径，但共同目标都是学到可迁移的视觉表示。
    \item 评价自监督方法时，最终还是要看下游任务表现，而不是只看预文本损失。
\end{enumerate}
\end{importantbox}

\subsection{下一讲预告}

\begin{figure}[H]
\centering
\includegraphics[width=0.95\textwidth]{slides-images/slide-111.jpg}
\caption{下一讲：Generative Models\protect\footnotemark}
\end{figure}
\footnotetext{Slides 第111页。}

下一讲我们会转向生成模型，讨论如何从“学表示”进一步走向“生成内容”。

\end{document}
